\documentclass{beamer}

\usetheme{Madrid}


\title{Measure Theoretic Probability}
\author{Edgar Yepez}
\date{December 3, 2027}

\begin{document}

\begin{frame}
  \titlepage
\end{frame}

\begin{frame}{Abstract}
  Introductory Probability courses typically bypass measure theoretic Probability in preference of caculus based methods. In this presentation we will build probability using measure thoery, intorducing probability space, sigma algebra, and probability measure. Using this framework we will define the random variable and expand the idea to stochastic processes and brownian motion.
\end{frame}

\begin{frame}{Outline}
  \tableofcontents
\end{frame}

% ==========================
\section{Introduction}
% ==========================
\begin{frame}{Motivation}
  Main motivation was to get a deeper understanding of stochastic calculus for its applications on the black scholes merton equation. 
\end{frame}


% ==========================
\section{Measure Theory}
% ==========================
\begin{frame}{What is a sigma-algebra}
  A sigma-algebra $\mathcal{F}$ is a collection of subsets of X such that 

  1. The empty set and the whole set are allways contained in the sigma algebra. 
  \[
    \varnothing, X \in \mathcal{F}
  \]
  2.If A is in the sigma algebra, then A complement is in the sigma algebra
  \[
    A \in \mathcal{F} \implies A^C \in \mathcal{F}
  \]
  3.If a countable collection of sets are in the sigma algebra then their untion is in the sigma-algebra. 
  \[
    \{A\}_{i=0}^{\infty} \subseteq \mathcal{F} \implies \bigcup_{i=0}^{\infty}A_i \in \mathcal{F}
  \]
\end{frame}

\begin{frame}{Sigma-algebra Examples}
  Let X = {1,2,3,4,5}, what is the sigma algebra generated by 
  $A_1 = \{1,2\}$, $A_2 = \{3,4\}$
  % answer: we know from the first axiom, empty set and X, 
  % since A_1 in our sigma algebra then its complemetn msut also be in  it  A_1^C = {3,4,5}, 
  % A_2^C = {1,2,5} 
  % unions and intersections {1,2,3,4} interesction of complenents {5} 
  \[
  \]

\end{frame}

\begin{frame}{What is a measure?}
  Given a set X with a sigma-alebra $\mathcal{F}$, then a $\textbf{measure}$ on $\mathcal{F}$ is definied as  function $\mu: \mathcal{F} \to [0,\infty]$ such that follows the following axioms:

   1. $\text{If } A\in \Sigma \text{ then its measure is non negative  }$
   \[
     0  \leq \mu(A) \leq \infty
   \]
   2. Measure of the empty set is 0 
   \[
      \mu(\varnothing) = 0 
   \]
   3. Given a sequence of pairwise disjoint sets $\{A\}_{i=0}^{\infty}$, that is for all i,j $A_i \cap A_j = \varnothing$, then 
  \[
      &\quad \mu(\bigcup_{i=0}^{\infty}(A)) = \bigcup_{i=0}^{\infty}\mu(A) 
  \]
\end{frame}



\begin{frame}{Lesbegue Measure}

\end{frame}

% ==========================
\section{Probability Theory}
% ==========================
\begin{frame}{Probability Theory}
  Probability Measure,

  Sigma algebra

  Probability Space 

  


\end{frame}

\begin{frame}{Random Variable}
 A random varible is a function $X: \Omega \to \mathbb{R}$

 Talk about expectation, independence , conditional probability

\end{frame}




% ==========================
\section{Stochastic Process}
% ==========================
\begin{frame}{Stochastic Process}
  A stockastic process is a sequence of random variables 

  \[
    \{X_t\}_{t=0}^T
  \]

  Expectation


  Variance


\end{frame}


\begin{frame}{Brownian Motion}
  
  \textbf{Definition}

  1. Brownian motion starts at 0 
  $W_0$ = 0
  2. Normably distributed mean 0 , varaince  t- s
  3. Independent increments 
\end{frame}

\begin{frame}{Example}
  Let
  \[
  Y_i =
  \begin{cases}
  1, & \text{with probability } \frac12,\\
  -1, & \text{with probability } \frac12.
  \end{cases}
  \]

  \[
    W_0=0,\qquad W_n=\sum_{i=1}^{n}\Delta x Y_i,\qquad \Delta x = \sqrt{\frac{t-s}{n}}
  \]
  
  Then $W_t - W_s = \sum_{i=s+1}^{t} \Delta x Y_i$.
  Since
  \[
    \mathbb{E}[\Delta x Y_i]  = \Delta x \mathbb{E}[Y_i] = \Delta x (\frac{1}{2} \cdot 1 + \frac{1}{2} \cdot (-1)) = 0 
  \]
  \[
    Var(\Delta x Y_i) = \Delta x^2 Var(Y_i) = \Delta x^2 (E[Y_i^2] - E[Y_i]^2)
  \]
\end{frame} 

\begin{frame}{Brownain Motion Illustration}
\end{frame}






\end{document}
